% ==== P4 draft F3 — §15.1 대체 블록 ====
\subsection{왜 흑연에는 없고 LCO 에는 있는가 --- MIT 의 물리}\label{sec:lco-why}
삽입 반응의 전체 엔트로피 변화는 ``리튬 자리 배열''(config)$\cdot$``격자 진동''(vib)$\cdot$``전자 준위 점유''
(electronic) 세 자유도의 합이며, 흑연에서는 셋째 몫이 작다 --- 흑연$\cdot$$\mathrm{LiC_6}$ 모두 전도성이 좋아
충방전 동안 $g(E_F)$ 가 크게 바뀌지 않아 전자 분포의 엔트로피 \emph{변화}가 거의 없다. LCO 는 다르다:
$x{=}1$($\mathrm{LiCoO_2}$)은 Co$^{3+}$($t_{2g}^6$ low-spin, 닫힌 껍질)라 \emph{전기 절연체}($g(E_F)\!\approx\!0$)
이며, 그 절연의 물리 부류는 단일입자 띠 그림 안에서 닫히는 \emph{밴드 절연체}(band insulator)다 --- 여섯
$3d$ 전자가 $t_{2g}$ 띠를 가득 채우고 그 위 $e_g$ 띠는 비어, Fermi 준위가 가득 찬 띠와 빈 띠 사이의
\emph{띠간극}(band gap)에 놓인다. 탈리튬화로 $x$ 가 $\sim$0.94 아래로 내려가면 Co$^{4+}$ 정공이 생겨
$t_{2g}$ 띠에 전도 전자가 열려 \emph{금속}이 된다($g(E_F)\!\to\!$유한)\cite{menetrier1999,motohashi2009}.
이 MIT 는 $x\approx0.75$--$0.94$ 의 2상역에서 일어나며(\S\ref{sec:lco-hys} 의 T1) 그 구간에서만 전자 분포가
켜지므로 전자 엔트로피 \emph{변화}도 그 구간에 국소한다 --- 곧 ``전자항이 켜지는 게이트''는 MIT 의 직접
결과다.

다만 이 켜짐의 \emph{꼴}은 밴드 그림 하나로는 닫히지 않는다: 단순 밴드 채움 논리라면 가득 찬 $t_{2g}$ 띠
꼭대기에 정공이 생기는 즉시 --- 임의로 작은 정공 농도에서 --- 금속이 되어야 하므로, $x{=}1$ 바로 아래에서
매끄럽게 시작되는 연속 전환을 예상하게 된다. 그러나 실측 MIT 는 위 조성 창의 2상 공존을 동반한 \emph{1차
전이}로 관측된다\cite{menetrier1999,reimers1992,motohashi2009} --- 정공이 왜 곧바로 퍼진 전도 상태가 되지
못하고 유한 조성 창에서 한꺼번에 풀리는가 하는 물음이 남는다. 이 물음의 배경과 해소(절연체의 두 부류와
금속--절연체 전이의 일반 물리, 그리고 실측 1차 MIT 를 설명하는 불순물 기작)는 별도의 배경 박스가
자족적으로 정리한다.

\begin{bgbox}
\textbf{금속--절연체 전이와 Mott 물리 ---} 고체가 절연체가 되는 길은 두 부류로 갈린다.
\emph{밴드 절연체}(band insulator)는 단일입자 띠 그림 안에서 닫히는 부류다 --- 전자들이 띠를 아래에서부터
채운 결과 Fermi 준위가 가득 찬 띠와 빈 띠 사이의 띠간극(band gap)에 놓이면, 전자 사이 상호작용을 부르지
않고도 절연이 설명된다. \emph{Mott 절연체}(Mott insulator)는 그 그림의 바깥이다 --- 띠가 부분적으로만 차
있어 띠 그림으로는 금속이어야 하나, 같은 자리에 전자 둘이 겹칠 때의 Coulomb 반발 $U$ 가 자리 사이
\emph{이동적분}(transfer integral) $t$ 를 누르면 전자가 이웃 자리로 건너가지 못하고 자리마다 하나씩
국소되어 절연한다 --- 띠 구조가 아니라 전자 \emph{상관}(correlation)이 만드는 절연이다\cite{mott1968,imada1998}.
금속상과 절연상을 가르는 \emph{금속--절연체 전이}(metal--insulator transition, MIT)는 이 셋 --- 띠
채움(도핑)$\cdot$띠폭(압력)$\cdot$상관 세기 --- 의 경쟁이 기우는 곳에서 일어난다\cite{imada1998}.
$\mathrm{LiCoO_2}$($x{=}1$)의 절연은 이 중 첫째 부류다 --- Co$^{3+}$($t_{2g}^6$ low-spin)의 가득 찬 $t_{2g}$
띠가 만드는 \emph{밴드 절연체}이고\cite{menetrier1999}, Mott 절연체와 공유하는 것은 ``절연''이라는 결과뿐이며
절연의 기원은 상관이 아니라 가득 찬 띠다. 밴드 절연체에 정공을 도핑하면 단순 띠 채움 논리로는 임의로 작은
정공 농도에서 곧바로 금속이 되어야 하나, 실측 $\mathrm{Li}_x\mathrm{CoO_2}$ 의 MIT 는 $x\approx0.75$--$0.94$
의 2상 공존을 동반한 \emph{1차 전이}로 관측된다\cite{menetrier1999,reimers1992,motohashi2009}.
Marianetti--Kotliar--Ceder 는 대형 초격자 DFT 계산으로 이 간극을 메우는 기작을 제시했다: 탈리튬화가 만드는
정공은 $t_{2g}$ 띠 전체로 곧장 퍼지는 대신 Li 빈자리에 속박된 \emph{불순물 준위}(impurity state)로 띠간극
안에 갇히고, 이 준위들이 이루는 불순물 띠(impurity band)가 임계 농도에서 한꺼번에 비국소화(delocalization)
하는 \emph{불순물 준위의 Mott 전이}가 1차 MIT 와 2상 공존을 설명한다\cite{marianetti2004}(계산은
$x\gtrsim0.95$ 절연$\cdot$$x\lesssim0.75$ 금속을 재현한다). 곧 상관 물리가 일하는 무대는 host $t_{2g}$ 띠가
아니라 탈리튬화가 심는 불순물 준위다 --- $x{=}1$ 끝점의 절연은 밴드 절연체인 채로 남고, Mott 물리는 불순물
띠의 소관이다. forward 모델은 이 미시 기작 전체를 Fermi 준위 상태밀도의 조성축 게이트 $g(E_F,x)$ 하나로
현상학화한다 --- 게이트 중심 $x_\mathrm{MIT}$ 는 실측 2상역의 중앙(전이가 일어나는 조성)이고, 폭
$\Delta x_\mathrm{MIT}$ 는 결함(산소 결손$\cdot$양이온 무질서)이 흩뜨리는 발현 조성의 분산이다(logistic
함수형과 세 초기값의 정당화는 \S\ref{sec:lco-gate} 가 닫는다). 미시 기작은 게이트가 \emph{왜} 그 조성에서
급격히 켜지는가를 답하고, 현상학 게이트는 그 켜짐이 조성축에서 \emph{어떻게} 진행하는가를 피팅 가능한
꼴로 적는다.
\end{bgbox}

켜진 전자 분포가 실제로 얼마만큼의 엔트로피를 담는가는 다음 소절(\S\ref{sec:lco-Se})이 Fermi--Dirac
분포에서 출발하는 유도로 닫는다.
% ==== 블록 끝 ====
% NOTE [Read Coverage]: AUTHOR_BRIEF.md 전문 / ch1_sec15_lcoelec.tex 전문(1–249행; §15.1 원문 + §15.4 게이트
%   정당화 역할 분담 확인) / ch1_sec13_lcohys.tex L120–164(T1 구조/전자 슬롯 분리 — 모순 없음 확인) /
%   V1020_STYLE_RUBRIC.md 전문 / ch1_sec02a_part0.tex L140–204(채택본 bgbox 선례 — 본문→박스 연결 관용
%   "별도의 배경 박스가 자족적으로 정리한다" 준용) / ch1_bib.tex(mott1968·imada1998·marianetti2004 키 실재 확인).
% NOTE [보존 확인 — 기존 §15.1 문장·인용 유지 목록]:
%   (1) 세 자유도 합(config·vib·electronic) + 흑연 셋째 몫 작음(전도성 좋아 g(E_F) 불변 → 변화 거의 없음): 원문 그대로.
%   (2) "LCO 는 다르다: x=1 은 Co3+(t2g6 low-spin, 닫힌 껍질)라 전기 절연체(g(E_F)≈0)": 원문 유지 +
%       밴드 절연체 부류·띠 그림(가득 찬 t2g·빈 e_g·띠간극) 정밀화만 덧붙임.
%   (3) "탈리튬화로 x~0.94 아래 → Co4+ 정공 → t2g 띠 전도 전자 → 금속(g→유한)" + \cite{menetrier1999,motohashi2009}: 원문 그대로.
%   (4) "이 MIT 는 x≈0.75–0.94 2상역(\S\ref{sec:lco-hys} 의 T1) ... 전자 엔트로피 변화도 그 구간에 국소 ---
%       게이트는 MIT 의 직접 결과다": 원문 그대로(마무리 다리 취지도 별도 문장으로 재전달).
% NOTE [자체검수 — 물리 가드]:
%   (G-a) x=1 을 Mott 절연체로 서술한 곳 0 — 본문·bgbox 모두 "밴드 절연체" 명시, bgbox 에 허용 대조 가드
%         ("공유하는 것은 '절연'뿐, 기원은 가득 찬 띠") 포함. PASS
%   (G-b) marianetti2004 과대 해석 없음 — "기작을 제시했다/설명한다" 수준, 상관 물리를 불순물 준위 소관으로
%         한정(host 띠와 대조), 유일 확정 기작 단정 없음. 수치는 브리프 검증 범위($x\gtrsim0.95$/$x\lesssim0.75$)만 사용. PASS
%   (G-c) 인용 6키(mott1968·imada1998·marianetti2004·menetrier1999·reimers1992·motohashi2009)만 사용,
%         수식·라벨 신설 0(산문+bgbox; 기존 \label{sec:lco-why} 유지). PASS
%   (G-d) 자족성 — bgbox 는 본문 지시대명사 없이 성립(\S\ref{sec:lco-gate} 는 위치 무관 해소), 본문은 bgbox
%         없이도 완결(물음 제기 후 박스 위임 문장만 남음). §15.4 와 역할 분담: bgbox=기작(왜), §15.4=게이트
%         함수형·세 초기값 정당화(어떻게) — 수치(0.85·0.05·13) 중복 배치 없음. PASS
